% --- 1. INFORMAZIONI GENERALI ---
\section{Informazioni Generali}
\begin{itemize}
	\item \textbf{Luogo/Piattaforma:} Server Discord
	\item \textbf{Data:} 2026-03-31
	\item \textbf{Orario di inizio:} 14:30
	\item \textbf{Orario di fine:} 14:50
	\item \textbf{Responsabile:} Sofia De Blasi
	\item \textbf{Scriba:} Sofia De Blasi
\end{itemize}

\vspace{0.5cm}
\textbf{Partecipanti presenti:}
\begin{table}[ht]
	\centering
	\renewcommand{\arraystretch}{1.2}
	\begin{tabularx}{\textwidth}{X l}
		\toprule
		\textbf{Nome e Cognome} & \textbf{Matricola} \\
		\midrule
		Sofia De Blasi & 2111014 \\
		Roberto Mariano Doroftei & 2111031 \\
		Filippo Idiometri & 2035617 \\
		Francesco Marcon & 2110990 \\
		Armando Moda Scarati & 2082864 \\
		Anton Ricardo Rupi & 2054304 \\
		\bottomrule
	\end{tabularx}
\end{table}

\vspace{0.5cm}
\textbf{Partecipanti assenti:}
\begin{table}[ht]
	\centering
	\renewcommand{\arraystretch}{1.2}
	\begin{tabularx}{\textwidth}{X l}
		\toprule
		\textbf{Nome e Cognome} & \textbf{Matricola} \\
		\midrule
		Nessuno & - \\
		\bottomrule
	\end{tabularx}
\end{table}

% --- 2. ORDINE DEL GIORNO ---
\section{Ordine del Giorno}
Gli argomenti previsti per la riunione odierna sono:
\begin{enumerate}
	\item Gestione chiamata con l'azienda Zucchetti del 2026-04-01.
	\item Suddivisione del lavoro e concentrazione su attività diverse.
	\item Pianificazione iniziale del progetto (Analisi dei Requisiti, Piano di Progetto, diagrammi di misurazione, gestione rotazione e sprint).
	\item Allineamento per tutti i membri sull'inizio dei casi d'uso (t4 - t5).
\end{enumerate}

\newpage
% --- 3. RESOCONTO ---
\section{Resoconto}

\subsection{Gestione chiamata aziendale Zucchetti}
Il gruppo ha discusso su come gestire l'imminente chiamata con l'azienda proponente Zucchetti, prevista per il giorno successivo, definendo i dettagli e i punti chiave da affrontare con il referente.

\subsection{Organizzazione inizio progetto e attività RTB}
La riunione si è focalizzata sull'organizzazione operativa delle prime fasi del progetto. Si è discusso di iniziare a strutturare i documenti chiave come l'Analisi dei Requisiti e il Piano di Progetto e la gestione della rotazione dei ruoli all'interno degli sprint. È stato inoltre concordato un allineamento generale per tutti i membri del gruppo, con l'obiettivo di mettersi in pari con le attività RTB (casi d'uso t4 e t5) entro venerdì.

\subsection{Gestione task e monitoraggio su GitHub Projects}
Per garantire un sistema standardizzato e un monitoraggio più efficace delle attività, è stato deciso di adottare la piattaforma GitHub Projects per la gestione di \textit{issue} e \textit{task}, integrando quanto già presente nel file Excel.

\subsection{Assegnazione prossimi verbali}
Il team ha concordato le assegnazioni per la stesura e la revisione dei prossimi verbali: il verbale dell'attuale riunione interna viene assegnato ad Armando (con Roberto come revisore); il verbale dell'incontro con Zucchetti sarà redatto da Sofia (con Filippo come revisore).

\newpage
% --- 4. DECISIONI ---
\section{Decisioni}

\begin{table}[ht]
	\centering
	\renewcommand{\arraystretch}{1.5}
	\begin{tabularx}{\textwidth}{c X}
		\toprule
		\textbf{Codice} & \textbf{Decisione} \\
		\midrule
		\textbf{VI-6.1} & Adozione di GitHub Projects come piattaforma principale per ordinare e monitorare le issue e task del progetto. \\
		\textbf{VI-6.2} & Allineamento generale sulle attività di inizio progetto (RTB, casi d'uso, Piano di Progetto). \\
		\textbf{VI-6.3} & Assegnazione formale dei ruoli per la stesura e revisione dei verbali della riunione odierna e dell'incontro con Zucchetti. \\
		\bottomrule
	\end{tabularx}
\end{table}

% --- 5. AZIONI DA SVOLGERE (TO-DO) ---
\section{Azioni da Svolgere (To-Do)}

\begin{table}[ht]
	\centering
	\renewcommand{\arraystretch}{1.5}
	\begin{tabularx}{\textwidth}{c c X l l}
		\toprule
		\textbf{Codice} & \textbf{Rif.} & \textbf{Descrizione Task} & \textbf{Assegnatario} & \textbf{Scadenza} \\
		\midrule
		\textbf{T-6.1} & VI-6.3 & Stesura del verbale interno del 2026-03-31 (Revisore: Roberto). & Armando & 2026-04-01 \\
		\textbf{T-6.2} & VI-6.3 & Stesura verbale esterno incontro Zucchetti del 2026-04-01 (Revisore: Filippo). & Sofia & 2026-04-02 \\
		\textbf{T-6.3} & VI-6.2 & Mettersi in pari con le attività RTB (t4 - t5 e inizio casi d'uso). & Tutti & 2026-04-03 \\
		\bottomrule
	\end{tabularx}
\end{table}

% --- 6. PROSSIMA RIUNIONE ---
\section{Prossime Riunioni}

\textbf{Incontro Esterno - Zucchetti}
\begin{itemize}
    \item \textbf{Data:} 2026-04-01
    \item \textbf{Ora:} 12:00 - 12:30
    \item \textbf{OdG preliminare:} Incontro esterno con l'azienda Zucchetti per chiedere aumento capienza gruppi.
\end{itemize}

\vspace{0.3cm}

\textbf{Riunione Interna \#7}
\begin{itemize}
    \item \textbf{Data:} 2026-04-03
    \item \textbf{Ora:} 14:30
    \item \textbf{OdG preliminare:} Allineamento e verifica dello stato di avanzamento delle attività RTB e dei casi d’uso. In caso di riscontro positivo da parte di Zucchetti circa l'aumento della capienza, si darà inizio alla pianificazione del progetto; in caso contrario, si avvierà la pianificazione per il contatto di un'azienda proponente alternativa.
\end{itemize}